\documentclass[11pt,a4paper]{article}
\usepackage[margin=2.5cm]{geometry}
\usepackage[T1]{fontenc}
\usepackage{lmodern}
\usepackage{amsmath,amssymb,amsthm}
\usepackage{booktabs}
\usepackage[colorlinks=true,linkcolor=blue,citecolor=blue,urlcolor=blue]{hyperref}
\newtheorem{result}{Result}
\theoremstyle{remark}
\newtheorem{remark}{Remark}
\numberwithin{equation}{section}

\title{Exact results for Kerr black holes: shadow areas and perimeters, capture of slow particles, and escape from near extremal horizons}
\author{Jacob Goodchild\thanks{Independent researcher. Contact: jacob.l.goodchild@gmail.com.
This work was carried out with substantial assistance from an AI system; see the
AI disclosure statement at the end of the paper.}}
\date{October 2026}

\begin{document}
\maketitle

\begin{abstract}
We give exact results for several observables of Kerr black holes that have so far been computed
numerically, by fits, or only at isolated parameter values. (i) The area of the shadow seen by an equatorial
observer, for every spin, in complete elliptic integrals, together with its small-spin series and its
near-extremal expansion $\mathcal A=16\pi+15\sqrt3+(3\pi/\sqrt2)\sqrt{1-a}+\dots$; the shadow area of an
extremal black hole at every inclination; and a single short integral for arbitrary spin and inclination.
(ii) The perimeter of the shadow, equal to $18\sqrt3\,M$ for a maximally spinning hole seen edge-on.
(iii) The capture cross-section of a Kerr black hole for slow particles arriving perpendicular to the spin,
$\sigma=[7\pi+\pi\sqrt{1-a^2}+16\sqrt{1+a}\,E(2a/(1+a))]M^2/v^2$, the corresponding angular-momentum accretion in closed form,
and the isotropically averaged cross-section and spin-down at maximal spin in elementary closed form.
(iv) Escape probabilities for photons emitted isotropically just outside the horizon of an extremal
Kerr--Newman black hole at any latitude, $P=\tfrac12-\arcsin(k)/(2\pi)-k/4$, and from the innermost stable circular
orbit at any charge, together with escaping energy fractions and escape probabilities for massive particles.
Known special cases are recovered and credited. All results were checked against independent high-precision
numerical computations.
\end{abstract}


\section{Introduction}

The Kerr metric is the unique stationary, asymptotically flat vacuum black hole of general relativity,
and many of its observable properties reduce to integrals over the critical (marginally trapped)
geodesics. For the shadow---the region of the observer's sky bounded by the image of the unstable
spherical photon orbits~\cite{Bardeen1973,Chandrasekhar1983}---the boundary curve is known in closed
parametric form, but derived quantities such as its \emph{area} and \emph{perimeter} have been computed
numerically or by interpolation; for instance, Cunha, Herdeiro and Radu remark that they ``were not able to
find an exact expression for the Kerr shadow areal radius in the generic case''~\cite{CunhaHerdeiroRadu2019}.
Likewise the capture cross-section of a Kerr black hole for slow (marginally bound) particles, relevant for
black holes moving through dark matter, has been obtained numerically or by fits~\cite{Young1976,Will2012,Karydas2026,MachMomenniaSarbach2026}, and the
probability that light or matter escapes from just outside the horizon of an extremal black hole has been
computed in closed form only at isolated points~\cite{YanGuoChen2021,GatesHadarLupsasca2021}.

In this paper we give exact formulas for these quantities:
\begin{itemize}
\item the area of the Kerr shadow seen edge-on, for every spin, in complete elliptic integrals,
with its small-spin and near-extremal expansions (Section~\ref{sec:edge});
\item the shadow area of an extremal Kerr black hole at every inclination, and a single-integral
representation for arbitrary spin and inclination (Sections~\ref{sec:ext}--\ref{sec:master});
\item the perimeter of the shadow (Section~\ref{sec:perim});
\item the edge-on capture cross-section for slow particles, the associated angular-momentum accretion,
and the isotropically averaged cross-section and spin-down at maximal spin (Section~\ref{sec:capture});
\item escape probabilities and escaping energy fractions for photons and massive particles emitted
isotropically near the horizon of an extremal Kerr(--Newman) black hole, at any latitude and for emitters
at rest or on the innermost stable circular orbit (Section~\ref{sec:escape}).
\end{itemize}
Results that were known before are identified as such in the text, with references. Units are
$G=c=M=1$ throughout; $a=J/M^2\in[0,1]$.

\section{Shadow boundary}\label{sec:setup}
For a distant observer at inclination $\vartheta$ to the spin axis ($s=\sin\vartheta$, $c=\cos\vartheta$),
the shadow boundary consists of the images of the spherical photon orbits of radius $r$, with
impact parameters~\cite{Bardeen1973,Chandrasekhar1983}
\begin{equation}
\xi(r)=\frac{r^2(3-r)-a^2(r+1)}{a(r-1)},\qquad
\eta(r)=\frac{r^3\big(4a^2-r(r-3)^2\big)}{a^2(r-1)^2},
\end{equation}
and celestial coordinates $\alpha=-\xi/s$, $\beta=\pm\sqrt{\eta+a^2c^2-\xi^2c^2/s^2}$. The area is
$\mathcal A=2\int\beta\,|d\alpha|$ over the range of $r$ where $\beta^2\ge0$. A key simplification is
\begin{equation}
\frac{d\xi}{dr}=-\frac{2\big[(r-1)^3+1-a^2\big]}{a(r-1)^2}.
\label{eq:dxi}
\end{equation}
The prograde and retrograde circular photon orbits are
$r_{p}=2\big[1+\cos\big(\tfrac23\arccos(-a)\big)\big]$, $r_{r}=2\big[1+\cos\big(\tfrac23\arccos a\big)\big]$; the third root of
$r^3-6r^2+9r-4a^2$ is $r_0=6-r_p-r_r\in(0,1)$.

\section{Edge-on shadow area for every spin}\label{sec:edge}
\begin{result}\label{res:edge}
For $\vartheta=\pi/2$ and $0<a<1$,
\begin{equation}
\mathcal A(a)=3\int_{r_p}^{r_r}\frac{(r-1)(5r-2)+2(1-a^2)/(r-1)}{\sqrt{r\,(4a^2-r(r-3)^2)}}\,dr .
\label{eq:edgeint}
\end{equation}
In terms of complete elliptic integrals $K=K(m)$, $E=E(m)$, $\Pi(n)=\Pi(n|m)$ (parameter convention),
with
\[
m=\frac{(r_r-r_p)r_0}{(r_r-r_0)r_p},\quad n=\frac{r_r-r_p}{r_r-r_0},\quad n'=n\frac{r_0-1}{r_p-1},\quad
g=\frac{2}{\sqrt{(r_r-r_0)r_p}},
\]
\[
V=\frac{nE+(m-n)K+(2nm+2n-n^2-3m)\Pi(n)}{2(n-1)(m-n)},
\]
\[
I_0=gK,\quad I_1=g\big[r_0K+(r_p-r_0)\Pi(n)\big],\quad
I_2=g\big[r_0^2K+2r_0(r_p-r_0)\Pi(n)+(r_p-r_0)^2V\big],
\]
\[
J=g\Big[\frac{K}{r_0-1}+\frac{r_0-r_p}{(r_p-1)(r_0-1)}\Pi(n')\Big],
\]
\begin{equation}
\mathcal A(a)=6I_0-21I_1+15I_2+6(1-a^2)J .
\label{eq:edgeclosed}
\end{equation}
\end{result}
\emph{Derivation.} Using \eqref{eq:dxi}, the area integrand is
$(4/a^2)\,r\sqrt{W}\,[1+(1-a^2)/(r-1)^3]$ with $W=r(4a^2-r(r-3)^2)$. Subtracting an exact derivative
$\frac{d}{dr}[S(r)\sqrt W]$ that vanishes at both endpoints (found by Hermite reduction with \textsc{SymPy})
leaves \eqref{eq:edgeint}; remarkably the polynomial part of the reduced integrand does not depend on $a$.
The substitution $\mathrm{sn}^2u=(r_r-r_0)(r-r_p)/[(r_r-r_p)(r-r_0)]$ gives \eqref{eq:edgeclosed}.

\paragraph{Expansions.} The small-spin expansion has rational coefficients (identified from 250-digit
evaluations of \eqref{eq:edgeclosed}; all denominators are of the form $2^p3^q$):
\begin{equation}
\frac{\mathcal A}{27\pi}=1-\frac{a^2}{18}-\frac{a^4}{72}-\frac{2059\,a^6}{314928}-\frac{88175\,a^8}{22674816}
-\frac{355999\,a^{10}}{136048896}-\dots
\end{equation}
The $a^2$ term follows from the known order-$a^2$ elliptical shape of the shadow~\cite{Bozza2006,GrallaLupsasca2020shape}
and is stated explicitly for the area in~\cite{Hassanabadi2026}; the $a^4$ term is contained in the perturbative results of
Kobialko and Gal'tsov~\cite{KobialkoGaltsov2026}. The higher coefficients appear to be new. At maximal spin $\mathcal A(1)=16\pi+15\sqrt3$, a value
obtained previously in~\cite{CunhaHerdeiroRadu2019,FengLu2020}. Near extremality, with $\epsilon=1-a$,
\begin{equation}
\mathcal A=16\pi+15\sqrt3+\frac{3\pi}{\sqrt2}\,\epsilon^{1/2}
+\epsilon\Big[\frac{8}{\sqrt3}\ln\frac{27}{2\epsilon}-6\sqrt3\Big]+o(\epsilon).
\label{eq:nearext}
\end{equation}
The $\epsilon^{1/2}$ term follows from matched asymptotics: as $a\to1$ the roots $r_0,r_p$ merge at $r=1$,
$r_{p,0}\simeq1\pm\delta$ with $\delta=\sqrt{8\epsilon/3}$; in the inner region the pole term
$6(1-a^2)/(r-1)\simeq12\epsilon/(r-1)$ dominates, and
$\int_1^\infty dx/(x\sqrt{x^2-1})=\pi/2$ gives $(12\epsilon/\sqrt3\delta)(\pi/2)=3\pi\epsilon^{1/2}/\sqrt2$.
The $\epsilon$ and $\epsilon\ln\epsilon$ coefficients were identified numerically (evaluating \eqref{eq:edgeclosed} with 60--150 digits
for $\epsilon$ between $10^{-12}$ and $10^{-80}$, followed by an integer-relation search). The square-root onset means that the edge-on shadow area
changes rapidly as the spin drops below its maximum.

\section{Extremal Kerr at every inclination}\label{sec:ext}
For $a=1$, $\xi=1+2r-r^2$ and $\eta=r^3(4-r)$. The quartic that defines $\beta^2$ factorises:
\begin{equation}
s^2\Big[r^3(4-r)+c^2-\frac{c^2}{s^2}(r^2-2r-1)^2\Big]=-\big(r^2-2(1+s)r-c^2\big)\big(r^2-2(1-s)r-c^2\big),
\end{equation}
with the elementary roots
$r_A=1-s-\sqrt{2-2s}<r_B=1+s-\sqrt{2+2s}<r_C=1-s+\sqrt{2-2s}<r_D=1+s+\sqrt{2+2s}$.
(Geometrically, Gralla and Lupsasca showed that the extremal critical curve is the convex hull of a
Cartesian oval, completed by the straight ``NHEK line'' when $\sin\vartheta>\sqrt3-1$~\cite{GrallaLupsasca2020shape};
the factorisation above is the corresponding statement in terms of the orbit radius.)
\begin{result}\label{res:ext}
With $\mathcal R(r)=(r-r_A)(r-r_B)(r-r_C)(r_D-r)$ and $\vartheta_c=\arcsin(\sqrt3-1)\approx47.0586^\circ$,
\begin{equation}
\mathcal A_{\rm ext}(\vartheta)=4\int_{r_{\rm lo}}^{r_D}\frac{3r^2-\cos^2\vartheta}{\sqrt{\mathcal R(r)}}\,dr+B(\vartheta),
\label{eq:extint}
\end{equation}
where $r_{\rm lo}=r_C$, $B=0$ for $\vartheta\le\vartheta_c$, and $r_{\rm lo}=1$,
$B=(2+s^2)\sqrt{3s^2+s^2c^2-4c^2}/s^2$ for $\vartheta>\vartheta_c$ (the NHEK-line contribution). For
$\vartheta\le\vartheta_c$, with $(e_1,e_2,e_3,e_4)=(r_D,r_C,r_B,r_A)$,
$m=\frac{(e_1-e_2)(e_3-e_4)}{(e_1-e_3)(e_2-e_4)}$,
$n=\frac{e_1-e_2}{e_1-e_3}$, $g=2/\sqrt{(e_1-e_3)(e_2-e_4)}$, $V$ as in
Result~\ref{res:edge},
\begin{equation}
\mathcal A_{\rm ext}=4\big(3I_2-c^2I_0\big),\quad I_0=gK,\quad
I_2=g\big[e_3^2K+2e_3(e_2-e_3)\Pi(n)+(e_2-e_3)^2V\big].
\end{equation}
\end{result}
The reduction uses, with $x=r-1$ and $P(x)\equiv\mathcal R(1+x)$, the identity $x\sqrt P=s^2(3x^2+6x+s^2+2)/\sqrt P+\frac{d}{dx}\big[(x^2-s^2-2)\sqrt P/4\big]$.
The limits are the known values $\pi(12+8\sqrt2)$ at $\vartheta=0$~\cite{FengLu2020} and $16\pi+15\sqrt3$ at
$\vartheta=\pi/2$. Near the axis,
\begin{equation}
\mathcal A_{\rm ext}=\pi(12+8\sqrt2)+\sqrt2\,\pi\Big[\frac{s^2}{2}+\frac{s^4}{128}-\frac{7s^6}{2048}-\frac{1735s^8}{524288}-\dots\Big]
\end{equation}
(coefficients identified by integer-relation search). For $a=1-\epsilon$ the area grows as
$\mathcal A_{\rm ext}(\vartheta)+c(\vartheta)\sqrt\epsilon$ with
\begin{equation}
c(\vartheta)=\frac{\pi\,(3-6\cos^2\vartheta-\cos^4\vartheta)}{\sqrt2\,\sin^3\vartheta}\quad(\vartheta>\vartheta_c),\qquad c(\vartheta)=0\quad(\vartheta<\vartheta_c),
\end{equation}
derived as in Section~\ref{sec:edge} from the master formula below; equivalently $c=\pi\beta_m^2/(\sqrt2\sin\vartheta)$
with $\beta_m$ the half-length of the NHEK line.
Sample values: $\mathcal A_{\rm ext}(17^\circ)=73.43231016343180$, $\mathcal A_{\rm ext}(45^\circ)=74.35818308197668$,
$\mathcal A_{\rm ext}(60^\circ)=75.08875773057813$.

\section{Single-integral formula for any spin and inclination}\label{sec:master}
For general $(a,\vartheta)$ the boundary involves a sextic, and no reduction to complete elliptic integrals is
expected (the corresponding curve has genus two). Nevertheless the area reduces to one short integral.
\begin{result}\label{res:master}
With $u=a^2\cos^2\vartheta$,
\[
N(r)=-r^6+6r^5-(9+2u)r^4+4a^2r^3+6ur^2-4a^2ur-u^2(r-1)^2
\]
(so that $N=\beta^2a^2(r-1)^2\sin^2\vartheta$) and $r_1<r_2$ the roots of $N$ with $r>1$ that bound the shadow,
\begin{multline}
\mathcal A(a,\vartheta)=\frac{1}{u-1}\int_{r_1}^{r_2}\Big[2(u^2-5a^2u+6u+3a^2-3)-2(2u^2+a^2u-3u-3a^2+3)r
+2(u^2-6u-3)r^2\\+4(u+3)r^3+2(u-3)r^4-\frac{2(a^2-1)(u^2+6u-3)}{r-1}\Big]\frac{dr}{\sqrt{N(r)}} .
\label{eq:master}
\end{multline}
\end{result}
At $\vartheta=\pi/2$ this reduces to \eqref{eq:edgeint}. The integrand has only inverse-square-root endpoint
singularities, so tanh--sinh quadrature converges to 15--30 digits at negligible cost.

\section{Perimeter of the shadow}\label{sec:perim}
\begin{result}\label{res:perim}
Along the critical curve the arc-length element simplifies to
$(d\ell/dr)^2=16r(r^2-a^2c^2)[(r-1)^3+1-a^2]^2/((r-1)^4N(r))$, so for any $(a,\vartheta)$
\begin{equation}
P(a,\vartheta)=8\int_{r_1}^{r_2}\frac{\big[(r-1)^3+1-a^2\big]\sqrt{r(r^2-a^2\cos^2\vartheta)}}{(r-1)^2\sqrt{N(r)}}\,dr .
\end{equation}
Edge-on only a cubic remains,
$P(a)=8\int_{r_p}^{r_r}\big[(r-1)+(1-a^2)/(r-1)^2\big]\,dr/\sqrt{4a^2-r(r-3)^2}$, and in closed form, with
$m=(r_r-r_p)/(r_r-r_0)$, $n=(r_r-r_p)/(r_r-1)$ and $V$ as above,
\begin{equation}
P(a)=\frac{16}{\sqrt{r_r-r_0}}\Big[r_0K+(r_r-r_0)E-K+\frac{(1-a^2)V}{(r_r-1)^2}\Big].
\end{equation}
At maximal spin edge-on, $d\ell=4\,dr/\sqrt{4-r}$ on $1\le r\le4$, so the curved part (both halves) has length $16\sqrt3$ and the
NHEK line $2\sqrt3$:
\begin{equation}
P(1,\pi/2)=18\sqrt3\approx31.17691 .
\end{equation}
\end{result}
Wei, Liu and Mann~\cite{WeiLiuMann2019} reported the length $16\sqrt3$ at $a=1$; this is the length of the curved
part only. Including the straight NHEK segment, which is part of the critical curve, gives $18\sqrt3$, continuous
with their value $31.2636$ at $a=0.99$. For small spin (coefficients identified at 60 digits),
\begin{equation}
\frac{P}{6\sqrt3\pi}=1-\frac{(1+C)a^2}{36}-\Big(\frac{35}{5184}+\frac{35C}{7776}-\frac{23C^2}{15552}\Big)a^4-\dots,\qquad C=\cos^2\vartheta .
\end{equation}

\section{Capture of slow particles}\label{sec:capture}
For particles arriving from infinity with speed $v\ll1$ the capture boundary in the plane of specific
angular momenta is traced by the marginally bound ($E=1$) spherical orbits. Writing $r=y^2$,
\begin{equation}
L_z=-\frac{y^4-2y^3+a^2}{a(y-1)},\qquad Q=\frac{y^4\big(a^2-(y^2-2y)^2\big)}{a^2(y-1)^2},\qquad 1+\sqrt{1-a}\le y\le1+\sqrt{1+a},
\label{eq:mb}
\end{equation}
which is equivalent to the marginally bound Carter constant given by Teo~\cite{Teo2021}. In particular
$L_z^2+Q=(3y^4-4y^3+a^2)/(y-1)^2$. The impact parameters are $(L_z,\sqrt Q)/v$ and the cross-section is
$\sigma=(M/v)^2\mathcal S$ with $\mathcal S=2\int\sqrt Q\,|dL_z|$.

\begin{result}[Edge-on cross-section]\label{res:slow}
For particles arriving perpendicular to the spin axis,
\begin{equation}
\sigma=\frac{M^2}{v^2}\Big[7\pi+\pi\sqrt{1-a^2}+16\sqrt{1+a}\;E\Big(\frac{2a}{1+a}\Big)\Big],
\label{eq:slow}
\end{equation}
which gives $16\pi M^2/v^2$ at $a=0$ and $(7\pi+16\sqrt2)M^2/v^2$ at $a=1$.
\end{result}
With $u=y-1$ the square root becomes $\sqrt{a^2-(u^2-1)^2}$, an even quartic; Hermite reduction gives
$\mathcal S=\int_{\sqrt{1-a}}^{\sqrt{1+a}}\big(16u^2+14u+2(1-a^2)/u\big)du/\sqrt{a^2-(u^2-1)^2}$, whose odd terms are
elementary. The expansion $\mathcal S/(16\pi)=1-a^2/16-31a^4/2048-233a^6/32768-\dots$ agrees through $a^4$ with
the slow limits of the fitted coefficients of Karydas, Vicente and Bertone~\cite{Karydas2026}, and the $a^2$ term with
Mach, Momennia and Sarbach~\cite{MachMomenniaSarbach2026}.

\begin{result}[Angular momentum accreted, edge-on]\label{res:M1}
The first moment of the capture region, $M_1=\iint L_z\,dL_z\,d\sqrt Q$, is
\begin{equation}
M_1=-\frac{32\sqrt{1+a}\,\big[(a^2+2)E-2(1-a)K\big]+35\pi a^2+10\pi\big(1-\sqrt{1-a^2}\big)}{5a},
\end{equation}
with $K=K(2a/(1+a))$, $E=E(2a/(1+a))$, so a black hole moving through slow matter in its equatorial plane gains angular momentum per unit mass
$dJ/dM=M\,M_1/\mathcal S=-aM\big(1+\tfrac{3}{32}a^2+\tfrac{37}{1024}a^4+\dots\big)$, and
$M_1\to-(96\sqrt2+45\pi)/5$ as $a\to1$.
\end{result}
The $a$ and $a^3$ terms are consistent with the slow limits of the fits of~\cite{Karydas2026}.

\begin{result}[Isotropic flux]\label{res:iso}
For an isotropic flux of slow particles, the direction-averaged cross-section is
$\langle\sigma\rangle=(\pi/2)(M/v)^2\int_{-1}^{1}L_c(\mu)^2d\mu$, where $\mu=\cos i=L_z/L_c$ is the inclination of
the orbital angular momentum, giving
\begin{equation}
\langle\sigma\rangle=\frac{\pi}{2}\frac{M^2}{v^2}\left|\int_{1+\sqrt{1-a}}^{1+\sqrt{1+a}}
\frac{2y^3(3y^4-8y^3+6y^2-a^2)}{a(y-1)^2\sqrt{3y^4-4y^3+a^2}}\,dy\right| .
\end{equation}
At $a=1$ (where $3y^4-4y^3+1=(y-1)^2(3y^2+2y+1)$ and the boundary includes the straight piece $L_z=2$ for
$\mu\ge\sqrt{2/3}$~\cite{Will2012}) everything is elementary:
\begin{equation}
\langle\sigma\rangle_{a=1}=\pi\frac{M^2}{v^2}\Big[\frac{47}{9}+\frac{31\sqrt2}{6}+\frac{7\sqrt6}{9}
-\frac{\sqrt3}{54}\ln\frac{1+\sqrt2+\sqrt6}{1+\sqrt2}\Big]=45.27564306526540\,\frac{M^2}{v^2},
\end{equation}
i.e.\ $0.90073$ of the Schwarzschild value $16\pi M^2/v^2$ (Young reported the ratio $0.90$~\cite{Young1976}). The
mean axial angular momentum per unit captured mass is $\langle l_z\rangle=\frac23\int\mu L_c^3d\mu/\int L_c^2d\mu$, and at
$a=1$
\begin{equation}
\int_{-1}^{1}\mu L_c^3\,d\mu=-\frac{4876}{189}-\frac{790\sqrt2}{63}+\frac{628\sqrt6}{189}-\frac{2\sqrt3}{81}\ln\frac{1+\sqrt2+\sqrt6}{1+\sqrt2},
\end{equation}
\begin{equation}
\langle l_z\rangle=-0.8193268418025525\,M .
\end{equation}
\end{result}
The uniform measure over incoming directions and impact-plane azimuths induces the uniform measure on the
direction of the orbital angular momentum, which justifies the averages. The small-spin series
$\langle l_z\rangle=-(2a/3+a^3/15+\dots)M$ is contained in~\cite{MachMomenniaSarbach2026,MachMomenniaSarbach2026b}.
Young~\cite{Young1976} reported $-0.828$ for this quantity at $a=1$; our exact value differs by about 1\%. It was confirmed
by an independent brute-force computation ($-0.819328$, Section~\ref{sec:verify}); we have not been able to consult the full text
of~\cite{Young1976} to identify the origin of the difference (numerical accuracy or a different averaging convention).

\begin{remark}[Any speed, maximal spin]
For $a=1$ and particles of any energy $\mathcal E=(1-v^2)^{-1/2}$, $p^2=\mathcal E^2-1$, arriving edge-on, the critical orbits are
$L=\mathcal E(r+1)-(r-1)\sqrt{r(p^2r+1)}$, $Q=r^2\big[-p^2r^2+(2p^2-1)r+1+2\mathcal E\sqrt{r(p^2r+1)}\big]$ for $1\le r\le r_{\max}$, and the
substitution $r=1/(t^2-p^2)$ reduces the cross-section to an incomplete elliptic integral with the quartic
$t^4+2\mathcal Et^3-t^2-2\mathcal Ep^2t-p^4$. As $v$ increases, $\sigma v^2$ starts from $7\pi+16\sqrt2$ ($v\to0$), and $\sigma\to16\pi+15\sqrt3$ as $v\to1$.
\end{remark}

\section{Escape from near the horizon of an extremal black hole}\label{sec:escape}
Consider the extremal Kerr--Newman family $a^2+q^2=1$, with $q$ the charge (Kerr: $a=1$) and an emitter at $r\to1^+$ at polar angle
$\vartheta$, emitting isotropically in its rest frame. For extremal black holes a finite fraction of the
emission escapes to infinity even in this limit~\cite{Ogasawara2020,YanGuoChen2021,GatesHadarLupsasca2021}.

\begin{result}[Emitter at rest]\label{res:zamo}
For an emitter at rest in the zero-angular-momentum frame, let
$k=(1+a^2\cos^2\vartheta)/(2a\sin\vartheta)$. The escape probability of photons is
\begin{equation}
P=\frac12-\frac{\arcsin k}{2\pi}-\frac k4\quad(k<1),\qquad P=0\quad(k\ge1).
\label{eq:Pzamo}
\end{equation}
The escape cone is $\{n_\varphi>0,\ n_\vartheta^2<(k^{-2}-1)n_\varphi^2,\ (n_\varphi>k\ \text{or}\ n_r>0)\}$ in terms of the
emitter-frame unit direction $(n_r,n_\vartheta,n_\varphi)$.
\end{result}
On the equator of extremal Kerr ($k=1/2$) this gives $7/24$, the value of Yan, Guo and Chen~\cite{YanGuoChen2021}
(numerically $0.2916$ in~\cite{Ogasawara2020}). On the equator of extremal Kerr--Newman ($k=1/(2a)$) it reproduces the
numerical values of Ogasawara et al.~\cite{Ogasawara2020} (Table II there: $0.267$, $0.194$, $0.0911$ for
$a=0.9,0.7,0.55$; formula: $0.26736$, $0.19480$, $0.09112$) and their threshold $a>1/2$. For Kerr, $P>0$ iff
$\sin\vartheta>\sqrt3-1$, the latitude range singled out by Ogasawara and Igata~\cite{OgasawaraIgata2021}; the edge
$n_\vartheta^2=(k^{-2}-1)n_\varphi^2$ corresponds to the near-horizon limit of the exact escape-cone boundary of Baines and
Visser~\cite{BainesVisser2025}. (The 7/24 value refers to the zero-angular-momentum frame; cone boundaries
expressed in other tetrads give different solid angles.)

\emph{Derivation.} Put $r=1+\varepsilon x$ and $\lambda=\lambda_H-\varepsilon u/a$, $\lambda_H=(1+a^2)/a$. In the
zero-angular-momentum frame $n_\varphi=\Sigma_H/(a\sin\vartheta\,(u+2))$, $\Sigma_H=1+a^2\cos^2\vartheta$. The radial
potential is $(r-1)^2[(r+1)^2-\mathcal C]$ far from the horizon and $\varepsilon^2[(2x+u)^2-\mathcal Cx^2]$ near it,
with $\mathcal C=\eta+1/a^2=(u+2)^2(n_\vartheta^2+n_\varphi^2)$; escape requires $\mathcal C<4$, and either outward
emission or $u<0$ (inward photons then bounce at $x=|u|/(2-\sqrt{\mathcal C})$). The solid angle of the resulting cone
is $\pi-2\arcsin k+\pi(1-k)$.

\begin{result}[Emitter on the ISCO]\label{res:isco}
For extremal Kerr--Newman with $a\ge1/\sqrt2$ the innermost stable circular orbit lies on the horizon and moves
with speed $v=k=1/(2a)$ relative to the zero-angular-momentum frame. A photon emitter on it escapes with
probability
\begin{equation}
P_{\rm ISCO}=\frac12-\frac{\arcsin k}{2\pi}+\frac{k}{\pi\sqrt{1+k^2}}\arctan\sqrt{\frac{1+k^2}{1-k^2}},\qquad k=\frac1{2a}.
\label{eq:Pisco}
\end{equation}
\end{result}
At $a=1$ this is $5/12+\arctan\sqrt{5/3}/(\sqrt5\pi)=0.546455$, the value found by Gates, Hadar and
Lupsasca~\cite{GatesHadarLupsasca2021} (numerically $54.6\%$ in~\cite{Igata2020}). At $a=1/\sqrt2$ it equals
$3/8+\sqrt3/9=0.567450$. It follows from boosting the cone of Result~\ref{res:zamo}: in the emitter frame,
with $t$ the cosine of the angle to the orbital motion, all directions with $t>0$ escape, none with $t<-k$, and
for $-k<t<0$ a fraction $\frac1\pi\arcsin\frac{t+k}{k\sqrt{1-t^2}}$, so
$P=\frac12+\frac1{2\pi}\int_{-k}^0\arcsin\frac{t+k}{k\sqrt{1-t^2}}\,dt$, which integrates to \eqref{eq:Pisco}.

\begin{result}[Escaping energy fractions]\label{res:energy}
The fraction of the emitted energy (in the emitter frame) that reaches infinity is, for the emitter at rest on
the equator ($a\ge1/2$),
\begin{equation}
\mathcal F_{\rm rest}=\frac{4a^2-1}{32a}+\frac{\sqrt{4a^2-1}}{16}\qquad(\text{Kerr: }\tfrac3{32}+\tfrac{\sqrt3}{16}=0.202003),
\end{equation}
and for the ISCO emitter of extremal Kerr
\begin{equation}
\mathcal F_{\rm ISCO}=\frac{2\sqrt3}{9}+\frac{1}{20\pi}+\frac{13\arctan\sqrt{5/3}}{5\sqrt{15}\,\pi}=0.595642 .
\end{equation}
\end{result}
These follow by weighting the escape regions with the energy ratio $g=E_\infty/E_{\rm emit}$, which equals
$a\,n_\varphi$ for the emitter at rest and $a(t+k)/\sqrt{1-k^2}$ for the ISCO emitter. The ISCO value can also be
obtained by integrating the redshift-resolved flux of Gates, Hadar and Lupsasca~\cite{GatesHadarLupsasca2021}
(their Eq.~(102)); we did not find the integrated fraction stated there.

\begin{result}[Massive particles]\label{res:massive}
(i) For particles emitted isotropically with local speed $\beta$ by an emitter at rest on the equator of an extremal
Kerr--Newman black hole, escape is impossible for $\beta<1/\sqrt{1+a^2}$ (Kerr: $1/\sqrt2$), and
\begin{equation}
P=\frac12\Big(1-\frac{\sqrt{1-\beta^2}}{a\beta}\Big)\qquad\Big(\frac{1}{\sqrt{1+a^2}}\le\beta\le\frac{\sqrt3}{2}\Big);
\end{equation}
for Kerr, $P=1/8$ exactly at $\beta=0.8$. (ii) For particles launched isotropically with speed $\beta$ from the ISCO of
extremal Kerr, escape requires $\beta>(3\sqrt2-2)/7=0.320377$, and
\begin{equation}
P=\frac{2\beta+1-\sqrt{3(1-\beta^2)}}{4\beta}\qquad\Big(\frac{3\sqrt2-2}{7}\le\beta\le\sqrt{\tfrac23}\Big),
\end{equation}
with $P=1/4$ at $\beta=1/2$ and $P=1/2$ at $\beta=\sqrt{2/3}$.
\end{result}
These follow from the same near-horizon analysis with massive geodesics: escape requires the conserved energy per unit mass to satisfy $\mathcal E\ge1$, the absence of a
turning point in the far zone ($3\mathcal E^2-1>Q$ for Kerr, with $Q$ the Carter constant) and the near-zone bounce condition. Above the stated ranges the
probabilities are given by one-dimensional integrals.

\section{Verification}\label{sec:verify}
All results were checked against independent computations (scripts in the code repository):
\begin{itemize}
\item closed forms against direct quadrature of the original parametric boundary curves at 30--40 digits
(e.g.\ edge-on area at $a=0.5$: $83.5610541284615455783534950$ from both);
\item areas against polygon (shoelace) areas of the boundary curves (agreement $\sim10^{-7}$--$10^{-9}$, limited by
the number of vertices);
\item the extremal formula at $60^\circ$ and $80^\circ$ against the master formula at $a=1-10^{-12}$ including the
$c(\vartheta)\sqrt\epsilon$ correction;
\item slow-capture cross-sections against brute-force bisection with the exact $E=1$ radial potential
($a=0.5$: $49.4258$ vs $49.4259$; $a=0.9$: $46.8049$ vs $46.8048$), and the isotropic $a=1$ values against direct bisection
for $L_c(\mu)$ at 2000 inclinations ($\langle l_z\rangle=-0.819328$ vs $-0.819327$);
\item escape probabilities against ray-by-ray simulations: grids of up to $1.8\times10^5$ emission directions at
$r=1+10^{-7}$, exact constants of motion computed without cancellation, escape decided from the exact radial potential
including inward-launched rays (agreement within the grid resolution, $\sim10^{-3}$; e.g.\ Kerr equator $0.29172$
vs $7/24$, ISCO $a=0.8$: $0.55890$ vs $0.55890$, massive $\beta=0.8$: $0.12502$ vs $1/8$).
\end{itemize}


\section*{AI disclosure statement}
This work was carried out by the author with substantial assistance from an AI system, Claude
(Anthropic), used through the Claude Code software. Under the author's direction, the AI system
proposed the research directions, wrote and ran the computer-algebra and numerical code (Python
with SymPy and mpmath), performed the derivations and checks reported here, searched the
literature, and drafted the text of this paper. The author directed the work, set the requirements
(verification of every result, honest novelty assessment, no fabricated citations) and is
responsible for its publication. All results were checked numerically as described in the text;
nevertheless, readers should treat the novelty assessment, which is based on literature searches,
with appropriate caution.

\begin{sloppypar}
\section*{Code availability}
All code used to derive and check the results is available at
\url{https://github.com/JacobGoodchild/findformula2} (directory \texttt{code/}, and the
verification scripts \texttt{papers/paperB/verify\_paperB.py} and \texttt{papers/paperB/verify\_paperB\_escape.py}).
\end{sloppypar}

\begin{thebibliography}{99}
\bibitem{Bardeen1973} J.~M.~Bardeen, ``Timelike and null geodesics in the Kerr metric'', in \emph{Black Holes (Les Astres Occlus)}, Les Houches 1972, eds.\ C.~DeWitt and B.~S.~DeWitt (Gordon and Breach, New York, 1973), pp.\ 215--240.
\bibitem{Chandrasekhar1983} S.~Chandrasekhar, \emph{The Mathematical Theory of Black Holes} (Clarendon Press, Oxford, 1983).
\bibitem{CunhaHerdeiroRadu2019} P.~V.~P.~Cunha, C.~A.~R.~Herdeiro and E.~Radu, ``EHT constraint on the ultralight scalar hair of the M87 supermassive black hole'', Universe \textbf{5}, 220 (2019), arXiv:1909.08039.
\bibitem{FengLu2020} X.-H.~Feng and H.~L\"u, ``On the size of rotating black holes'', Eur.\ Phys.\ J.\ C \textbf{80}, 551 (2020), arXiv:1911.12368.
\bibitem{Bozza2006} V.~Bozza, F.~De~Luca and G.~Scarpetta, ``Kerr black hole lensing for generic observers in the strong deflection limit'', Phys.\ Rev.\ D \textbf{74}, 063001 (2006), arXiv:gr-qc/0604093.
\bibitem{GrallaLupsasca2020shape} S.~E.~Gralla and A.~Lupsasca, ``Observable shape of black hole photon rings'', Phys.\ Rev.\ D \textbf{102}, 124003 (2020), arXiv:2007.10336.
\bibitem{KobialkoGaltsov2026} K.~Kobialko and D.~Gal'tsov, ``Perturbation theory for gravitational shadows in Kerr-like spacetimes'', Phys.\ Rev.\ D \textbf{113}, 104007 (2026), arXiv:2512.24126.
\bibitem{Hassanabadi2026} H.~Hassanabadi, M.~R.~R.~Good, S.~Zare, O.~Luongo and F.~Kafikang, ``Optical and thermodynamic properties of Kerr-Bertotti-Robinson black holes'', arXiv:2607.11979 (2026).
\bibitem{WeiLiuMann2019} S.-W.~Wei, Y.-X.~Liu and R.~B.~Mann, ``Intrinsic curvature and topology of shadows in Kerr spacetime'', Phys.\ Rev.\ D \textbf{99}, 041303 (2019), arXiv:1811.00047.
\bibitem{Young1976} P.~J.~Young, ``Capture of particles from plunge orbits by a black hole'', Phys.\ Rev.\ D \textbf{14}, 3281 (1976).
\bibitem{Will2012} C.~M.~Will, ``Capture of non-relativistic particles in eccentric orbits by a Kerr black hole'', Class.\ Quantum Grav.\ \textbf{29}, 217001 (2012), arXiv:1208.3931.
\bibitem{Teo2021} E.~Teo, ``Spherical orbits around a Kerr black hole'', Gen.\ Rel.\ Grav.\ \textbf{53}, 10 (2021), arXiv:2007.04022.
\bibitem{Karydas2026} T.~K.~Karydas, R.~Vicente and G.~Bertone, ``Mass and spin coevolution of black holes inspiralling through dark matter'', Phys.\ Rev.\ D \textbf{113}, 043057 (2026), arXiv:2510.13604.
\bibitem{MachMomenniaSarbach2026} P.~Mach, M.~Momennia and O.~Sarbach, ``Accretion of a Vlasov gas by a Kerr black hole'', Phys.\ Rev.\ D \textbf{113}, 044068 (2026), arXiv:2508.04783.
\bibitem{MachMomenniaSarbach2026b} P.~Mach, M.~Momennia and O.~Sarbach, ``Bondi-type accretion onto a Kerr black hole in the kinetic regime'', Phys.\ Rev.\ D \textbf{113}, L041505 (2026), arXiv:2508.20189.
\bibitem{Ogasawara2020} K.~Ogasawara, T.~Igata, T.~Harada and U.~Miyamoto, ``Escape probability of a photon emitted near the black hole horizon'', Phys.\ Rev.\ D \textbf{101}, 044023 (2020), arXiv:1910.01528.
\bibitem{Igata2020} T.~Igata, K.~Nakashi and K.~Ogasawara, ``Observability of the innermost stable circular orbit in a near-extremal Kerr black hole'', Phys.\ Rev.\ D \textbf{101}, 044044 (2020), arXiv:1910.12682.
\bibitem{YanGuoChen2021} H.~Yan, M.~Guo and B.~Chen, ``Observability of zero-angular-momentum sources near Kerr black holes'', Eur.\ Phys.\ J.\ C \textbf{81}, 847 (2021), arXiv:2104.07889.
\bibitem{GatesHadarLupsasca2021} D.~E.~A.~Gates, S.~Hadar and A.~Lupsasca, ``Photon emission from circular equatorial Kerr orbiters'', Phys.\ Rev.\ D \textbf{103}, 044050 (2021), arXiv:2010.07330.
\bibitem{OgasawaraIgata2021} K.~Ogasawara and T.~Igata, ``Complete classification of photon escape in the Kerr black hole spacetime'', Phys.\ Rev.\ D \textbf{103}, 044029 (2021), arXiv:2011.04380.
\bibitem{BainesVisser2025} J.~Baines and M.~Visser, ``Exact analytic photon escape cones and silhouettes for extremal Kerr black holes'', Phys.\ Rev.\ D \textbf{111}, L021501 (2025), arXiv:2405.08875.

\end{thebibliography}
\end{document}
